% Options for packages loaded elsewhere
\PassOptionsToPackage{unicode}{hyperref}
\PassOptionsToPackage{hyphens}{url}
%
\documentclass[
]{article}
\usepackage{amsmath,amssymb}
\usepackage{iftex}
\ifPDFTeX
  \usepackage[T1]{fontenc}
  \usepackage[utf8]{inputenc}
  \usepackage{textcomp} % provide euro and other symbols
\else % if luatex or xetex
  \usepackage{unicode-math} % this also loads fontspec
  \defaultfontfeatures{Scale=MatchLowercase}
  \defaultfontfeatures[\rmfamily]{Ligatures=TeX,Scale=1}
\fi
\usepackage{lmodern}
\ifPDFTeX\else
  % xetex/luatex font selection
\fi
% Use upquote if available, for straight quotes in verbatim environments
\IfFileExists{upquote.sty}{\usepackage{upquote}}{}
\IfFileExists{microtype.sty}{% use microtype if available
  \usepackage[]{microtype}
  \UseMicrotypeSet[protrusion]{basicmath} % disable protrusion for tt fonts
}{}
\makeatletter
\@ifundefined{KOMAClassName}{% if non-KOMA class
  \IfFileExists{parskip.sty}{%
    \usepackage{parskip}
  }{% else
    \setlength{\parindent}{0pt}
    \setlength{\parskip}{6pt plus 2pt minus 1pt}}
}{% if KOMA class
  \KOMAoptions{parskip=half}}
\makeatother
\usepackage{xcolor}
\usepackage[margin=1in]{geometry}
\usepackage{graphicx}
\makeatletter
\def\maxwidth{\ifdim\Gin@nat@width>\linewidth\linewidth\else\Gin@nat@width\fi}
\def\maxheight{\ifdim\Gin@nat@height>\textheight\textheight\else\Gin@nat@height\fi}
\makeatother
% Scale images if necessary, so that they will not overflow the page
% margins by default, and it is still possible to overwrite the defaults
% using explicit options in \includegraphics[width, height, ...]{}
\setkeys{Gin}{width=\maxwidth,height=\maxheight,keepaspectratio}
% Set default figure placement to htbp
\makeatletter
\def\fps@figure{htbp}
\makeatother
\setlength{\emergencystretch}{3em} % prevent overfull lines
\providecommand{\tightlist}{%
  \setlength{\itemsep}{0pt}\setlength{\parskip}{0pt}}
\setcounter{secnumdepth}{-\maxdimen} % remove section numbering
\usepackage{booktabs}
\usepackage{longtable}
\usepackage{array}
\usepackage{multirow}
\usepackage{wrapfig}
\usepackage{float}
\usepackage{colortbl}
\usepackage{pdflscape}
\usepackage{tabu}
\usepackage{threeparttable}
\usepackage{threeparttablex}
\usepackage[normalem]{ulem}
\usepackage{makecell}
\usepackage{xcolor}
\ifLuaTeX
  \usepackage{selnolig}  % disable illegal ligatures
\fi
\usepackage{bookmark}
\IfFileExists{xurl.sty}{\usepackage{xurl}}{} % add URL line breaks if available
\urlstyle{same}
\hypersetup{
  pdftitle={Final Report},
  hidelinks,
  pdfcreator={LaTeX via pandoc}}

\title{Final Report}
\author{}
\date{\vspace{-2.5em}}

\begin{document}
\maketitle

\section{Introduction}\label{introduction}

It is 7:30 in the morning, and everything depends on getting to the
subway on time. You checked the schedule, left early just in case, and
felt confident. But now, as you step onto the platform, the signs are
blank, the train is missing, and the crowd is growing. For many in
Toronto, this is a familiar scene. Transit delays are part of daily
life. They cause lateness, missed appointments, and rising stress. Over
time, they weaken trust in the system and disrupt the flow of the entire
city.

The Toronto Transit Commission runs a complex network of subways, buses,
and streetcars. Delays affect not just individual commuters but also the
wider economy and city life. Better delay prediction can help TTC make
smarter decisions about planning, staffing, and live updates. It also
helps riders travel with more certainty. More broadly, accurate
forecasts benefit businesses, reduce traffic pressure, and support a
more efficient and resilient city.

This study asks whether these delays can be predicted more accurately by
using past delay patterns and signals from major external events. The
goal is to understand if yesterday's or last week's delays can help
forecast today's, and whether including a flag for major disruptions,
such as the COVID-19 stay-at-home period, improves those predictions.

Although COVID-19 is used in this study, it serves as an example rather
than the main focus. Major events like pandemics, severe weather, labor
strikes, or infrastructure failures can all affect transit patterns. By
using COVID as a case, this research explores whether models can
recognize and adapt to these kinds of disruptions.

The data used come from the TTC Open Data Portal, which provides public
access to transit delay records across subways, buses, and streetcars.
Covering the years 2014 to 2024, the dataset includes over 400,000 delay
events. To study the research questions, three types of models are
applied: an Autoregressive Integrated Moving Average model with external
variables, a Generalized Additive Model, and a Random Forest model. Some
models use only recent delay history, while others include a simple
signal to represent the impact of major disruptions. These approaches
help reveal whether daily transit delays follow regular patterns or are
strongly shaped by external events.

\section{Methods}\label{methods}

\subsection{Data}\label{data}

The data used in this study came from the Toronto Transit Commission
delay records available through the Toronto Open Data Portal. These
records cover transit delays from subway, streetcar, and bus services in
Toronto between 2014 and 2024. Subway and streetcar delay data were
downloaded directly from the online portal, while bus delay information
and a subway delay reference guide were downloaded separately as
spreadsheet files. The original data are freely accessible online
through the Toronto Open Data website.

Before analysis, the data required cleaning and standardization.
Originally, each delay was recorded with separate date and time fields.
These fields were combined into one single timestamp called
\texttt{DateTime}. This timestamp was initially in Coordinated Universal
Time (UTC) but was converted to the local Toronto time zone (Eastern
Standard Time). Delay durations, which were recorded differently across
datasets, were combined into one standardized variable named
\texttt{min\_delay}, measured in minutes. Delay incidents, initially
recorded with numeric codes in subway data, were replaced with clear
descriptions using the subway delay reference guide. Finally, any
records with missing or invalid information (such as negative delays)
were removed. After cleaning, the final dataset contained 400,845 delay
events: 207,014 from buses, 139,007 from streetcars, and 54,824 from
subways.The cleaned data can be found in the data folder, or retrieved
using the script in the data folder.

This report used a cleaned csv file, but all results can be fully
reproduced using the raw TTC delay data available online. The data/
folder contains the exact scripts used for data acquisition and cleaning
in order to get the cleaned csv file.

\subsection{Models}\label{models}

This study aims to predict the average daily delay across Toronto's
transit network. To calculate this, all delays for each day were summed
and then divided by the number of delay events that day. Predictions
used two kinds of information. First, historical delays from the day
before (``lag-1'') and from exactly one week ago (``lag-7'') were
included. These lagged delays help models understand if delays today
might be similar to recent delays or follow a weekly pattern. Second, an
indicator was used to mark a significant external event: the COVID-19
pandemic period. Days during the major COVID-19 stay-at-home orders
(March 1, 2020, to June 30, 2022) were labeled ``1'', and all other days
were labeled ``0''. This helps the models learn if delays were different
during COVID compared to normal times. To fairly test the models, the
last 14 days of data were reserved for evaluation, and all earlier data
were used for training. The seed was set to 370 for reproducibility.

A simple baseline forecast, seasonal naive, was first computed. This
basic method assumes the delay today will be exactly the same as the
delay from exactly one week ago. Mathematically, it is:
\(\widehat{y}_{t+h} = y_{t+h-7}\). Here, \(\widehat{y}_{t+h}\)
represents the forecasted delay \(h\) days from today. This simple
method provides a basic standard that more advanced models should aim to
outperform.

The first advanced model tested was an ARIMAX model. ARIMAX predicts
future delays by analyzing past delays and external events. It finds
patterns from past delays---daily and weekly---and checks whether delays
changed significantly during the COVID-19 period. The ARIMAX model is
expressed as:

\[
\text{Delay today} = (\text{Patterns from past delays}) + (\text{Effect of COVID event}) + (\text{Random error})
\]

This model calculates coefficients that measure how past delays and the
COVID period impact today's delays. A positive COVID coefficient means
delays were higher during COVID, and a negative one means delays were
lower. The best version of this model was chosen automatically using
AICc, which balances accuracy and simplicity.

To judge whether the pandemic indicator really helps, a second model was
fit that is identical in every way except it ignores the COVID flag. In
other words, the forecast is driven only by the series' own history; no
outside information is supplied. This plain ARTIMA model is written

\[
\text{Delay today} = (\text{Patterns from past delays}) + (\text{Random error})
\]

Next, a Generalized Additive Model was used. The GAM allows the delay
trend over time to be smooth and flexible, not necessarily following
straight lines. It helps identify if delays gradually increase or
decrease over certain periods and also examines how delays differed
before, during, and after COVID-19. The GAM model is written as:

\[
\text{Delay} = (\text{Average delay}) + (\text{Smooth trend over time}) + (\text{COVID period effect}) + (\text{Random error})
\]

The ``smooth trend'' is created using spline, which smoothly curves to
fit data without sudden jumps. GAM provides a measure called ``degrees
of freedom,'' which indicates how flexible or complex the curve is.
Higher values mean the model found more complex patterns. COVID-period
coefficients show whether delays during and after COVID differed
significantly from normal times.

A Random Forest model was also used to predict average daily transit
delays. Random Forest is an ensemble learning method that builds many
individual regression trees and then averages their predictions. Each
regression tree is a simple model that tries to predict today's average
delay based on a few input features: the delay from yesterday (lag1),
the delay from exactly one week ago (lag7), and an indicator for whether
the date falls during the COVID-19 period.

In a single tree, the data are repeatedly split into groups based on the
predictor variables. At each split, the goal is to divide the data into
two parts so that the delay values within each part are as similar as
possible. This splitting continues until a stopping rule is met, such as
reaching a minimum number of observations in a group. Each final group
predicts a single value, which is usually the average delay of the
points in that group. The Random Forest grows many of these trees, each
on a different random sample of the training data. It also randomly
selects a subset of predictors at each split to create diverse trees.
The final forecast is simply the average of the predictions from all the
trees. The idea behind the Random Forest model can be written as:

\[
\text{Delay today} = \frac{1}{T} \sum_{t=1}^T \text{Tree}_t(\text{lag1}, \text{lag7}, \text{COVID})
\]

where \(T\) is the total number of trees, and \(\text{Tree}_t(\cdot)\)
is the prediction made by the \(t\)-th regression tree using the given
predictors.

\paragraph{Criteria}\label{criteria}

To compare the performance of these models, three different evaluation
metrics were used: Root Mean Squared Error (RMSE), Mean Absolute Error
(MAE), and Mean Absolute Scaled Error (MASE). RMSE measures the average
size of the prediction errors, giving extra weight to large errors. MAE
measures the simple average size of errors, making it easy to interpret.
MASE compares the forecast errors from the models to the seasonal naive
baseline method; a MASE value below 1 means the model performed better
than the simple naive forecast. By comparing these metrics, it becomes
possible to determine whether routine delay patterns alone or
significant external events (such as the COVID-19 pandemic) provide more
valuable information when predicting daily transit delays.

\section{Results}\label{results}

The seasonal-naive benchmark assumes that the delay for any given day
will be identical to the delay recorded exactly one week earlier. When
this very simple rule was applied to the fourteen-day hold-out it
produced a root-mean-square error of about 7.04 minutes and a
mean-absolute error of roughly 5.38 minutes, so the typical forecast
missed the true value by just over five minutes. These numbers establish
a baseline that the more sophisticated methods are expected to beat.

\subsubsection{Time Series}\label{time-series}

The ARIMA model without the COVID flag estimates the daily-delay series
as

\[
(1 - 0.44\,B)\,(1-B)\,y_t
    \;=\;
    (1 - 1.35\,B + 0.36\,B^{2})\,
    (1 - 0.07\,B^{7} - 0.06\,B^{14})\,\varepsilon_t,
\]

where \(B\) is the one--day back-shift operator, \(y_t\) is the average
delay on day \(t\) and \(\varepsilon_t\) is random noise.All numeric
coefficients represent the trend. More specifically, \(0.44\) says
today's deviations depend moderately on yesterday's. The two terms
\(-1.35\) and \(0.36\) summaries short bursts of shocks. The terms
\(-0.07\) and \(-0.06\) capturs the TTC's regular seven-day rhythm. With
these terms the model achieves a hold-out RMSE of 3.78 minutes and a
MASE of 0.64, already much better than the simple seasonal-naive
benchmark.

The ARIMAX model with a COVID indicator keeps exactly the same lag
structure but adds a regression term for COVID, where \(x_t\) equals 1
on days between 1 March 2020 and 30 June 2022 and 0 otherwise:

\[
(1 - 0.75\,B)\,(1-B)\,y_t
   \;=\;
   (1 - 1.68\,B + 0.68\,B^{2})\,
   (1 - 0.58\,B^{7})\,\varepsilon_t
   \;-\;1.62\,x_t ,
\]

The negative coefficient \(-1.62\) means that, holding all routine
patterns fixed, the pandemic period lowers the expected daily delay by
about 1.6 minutes. None of the other numbers change significantly, which
implies the week-to-week trend of the series remains intact; only the
level shifts downward when \(x_t=1\). Adding this single term trims the
hold-out RMSE to 3.67 minutes and the MASE to 0.62, a modest two-percent
improvement.

Because both models share the same lag structure, their difference
isolates the effect of the COVID indicator. The small reduction in
forecasting error together with the unchanged AR and MA coefficients
supports the idea that COVID did not alter how delays evolve from day to
day or week to week; it simply moved the whole curve slightly lower.
Thus the long-run trend and seasonal pattern stayed the same, but the
pandemic introduced a temporary level drop of roughly one and a half
minutes in average daily delay.

\subsubsection{GAM}\label{gam}

The generalized additive model represents the daily delay as a smooth
curve of calendar time plus fixed adjustments for the three pandemic
phases. In equation form the fitted model is

\[
\hat y_i = 15.52 + f(\text{date}_i) - 1.80\,\mathbf 1_{\text{COVID}} - 1.60\,\mathbf 1_{\text{post}} ,
\]

where \(f(\cdot)\) is a cyclic cubic spline with about 8.7 effective
degrees of freedom. The spline captures gradual rises and falls that
repeat each year, while the two negative coefficients show that the
average delay dropped by roughly 1.8 minutes during the strict stay-home
interval and by about 1.6 minutes in the months that followed. Because
the same smooth curve is applied to all three phases, these results
indicate that the overall trend shape stayed intact and only its height
shifted downward while travel volumes were lower. Forecast accuracy
improved further, with an RMSE of 3.49 minutes, an MAE of 2.91 minutes
and a MASE of 0.57.

\subsubsection{Random Forest}\label{random-forest}

This model achieved the best performance out of all methods tested, with
a RMSE of 2.73 minutes, a MA of 2.22 minutes, and a MAS of 0.43 on the
14-day hold-out period. These results show that the Random Forest
provided highly accurate predictions, with errors noticeably smaller
than the baseline and the other models.

A variable importance analysis was also performed to understand which
predictors contributed the most to the Random Forest's accuracy. The
results are shown in the bar plot below. The plot displays the
importance scores for each variable, where a higher score means the
variable played a bigger role in making accurate predictions. The
COVID-19 indicator was the most important predictor with a score around
15.5, showing that the pandemic had a strong influence on delays even
after controlling for historical patterns. The previous day's delay
(lag1) had an importance score around 6.9, and the delay from one week
ago (lag7) had a score around 5.7. This means that while short-term and
weekly history were important, the COVID period itself had a noticeable
and lasting effect on daily delays.

\begin{center}\includegraphics[width=700px]{report_files/figure-latex/unnamed-chunk-3-1} \end{center}

\paragraph{Summary}\label{summary}

Taken together, these results paint a consistent picture. Every
modelling framework finds that the weekly rhythm and short-term momentum
of delays persisted through the pandemic, but the overall level of those
delays moved lower by roughly one-and-a-half to two minutes while
stay-home orders were in place. Including a simple binary flag for that
event therefore improves forecasts without needing to alter the lag
structure or the long-run trend shape. The hypothesis that exceptional
events enhance predictive accuracy is supported, yet the evidence also
shows that the dominant signal still comes from routine daily and weekly
patterns.

The chart compares each model's forecasts with the actual delays over
the last two weeks. The green line is the real data. The grey baseline
line, which simply repeats last week's value, often sits well above the
true delays, especially on the sharp dips. The two orange lines---plain
ARIMA and ARIMAX---follow almost the same smooth path; they are closer
to reality than the baseline but still miss the sudden drops. The blue
GAM curve stays almost flat, matching the overall level yet overlooking
the daily swings. The red Random Forest track stays nearest to the green
line, moving up and down at the right moments and therefore gives the
most accurate day-by-day predictions.

\begin{center}\includegraphics[width=700px]{report_files/figure-latex/unnamed-chunk-4-1} \end{center}

The table shows how well each model forecasts the next 14 days using
three error scores: RMSE, MAE and MASE. Smaller numbers mean better
accuracy. The baseline that simply repeats last week's delay is the
weakest, missing by about 7 minutes on RMSE and 5.4 minutes on MAE,
which sets MASE near 1.05. Adding an ARIMA structure that learns recent
patterns cuts those errors by more than half, and including the COVID
flag in ARIMAX trims them a little more. The GAM, which fits a smooth
calendar curve plus COVID shifts, lowers RMSE to about 3.5 minutes. The
random forest, which combines yesterday's delay, the delay a week ago
and the COVID flag across many decision trees, is the most accurate.
Overall, every model that blends recent history with the pandemic
indicator beats the weekly-repeat baseline, and the nonlinear random
forest delivers the biggest improvement.

\begin{longtable}[t]{lrrr}
\caption{\label{tab:unnamed-chunk-5}Model Performance Comparison (14-day Hold-Out Set)}\\
\toprule
Model & RMSE & MAE & MASE\\
\midrule
Baseline & 7.04 & 5.38 & 1.05\\
ARIMA & 3.78 & 3.26 & 0.64\\
ARIMAX & 3.67 & 3.18 & 0.62\\
GAM-COVID & 3.49 & 2.91 & 0.57\\
RF & 2.73 & 2.22 & 0.43\\
\bottomrule
\end{longtable}

\section{Conclusions}\label{conclusions}

Transit delays affect more than just people trying to get to work or
school. They shape how smoothly a city runs. This study shows that many
of these delays follow regular patterns. Delays from the day before or
the same day last week can help predict what will happen today. These
patterns make it possible to build models that give useful forecasts.

The study also finds that adding a simple signal for major events, like
the COVID-19 pandemic, can make predictions better. These events do not
change how delays behave from day to day, but they can shift the overall
level of delays. During the COVID stay-at-home period, for example, the
average delay dropped. This shows that outside events, even if rare, can
have a strong effect.

These findings matter because cities need to plan for both normal days
and times of disruption. If transit systems can see when delays are
likely to happen, they can respond faster. This helps with staffing,
service planning, and letting riders know what to expect. When transit
is more reliable, people trust it more and are more likely to use it.

There are also limits to what this study shows. The models used here
only looked at past delays and one simple event flag. They did not
include other factors like weather, traffic, or how many people were
riding each day. These might also affect delays and could be useful in
future studies. The models also do not explain why a disruption happened
or how long its effects lasted.

Future work could use more detailed data and test new kinds of models.
It could also explore how to use delay forecasts in real time to support
decisions made by transit staff. These tools could help make public
transportation more reliable not just in Toronto, but in other cities
too.

\end{document}
